% !TEX program = pdflatex
\documentclass[11pt,a4paper]{article}
\usepackage[utf8]{inputenc}
\usepackage[T1]{fontenc}
\usepackage[brazilian]{babel}
\usepackage{lmodern}
\usepackage[margin=2.0cm,top=2.0cm,bottom=2.2cm]{geometry}
\usepackage{xcolor}
\usepackage{enumitem}
\usepackage{fancyhdr}
\usepackage{tcolorbox}
\usepackage{amsmath,amssymb}
\usepackage{microtype}
\usepackage[hidelinks]{hyperref}

\definecolor{verde}{HTML}{24574D}
\definecolor{azul}{HTML}{1769A5}
\definecolor{grafite}{HTML}{253A40}
\definecolor{cinza}{HTML}{7F8C8D}
\definecolor{cinzaclaro}{HTML}{F8F9FA}
\definecolor{alerta}{HTML}{C0392B}

\newcommand{\slidenum}{1}
\pagestyle{fancy}
\fancyhf{}
\lhead{\color{cinza}\footnotesize\bfseries Tópicos em Microeconomia Aplicada --- 16 de Novembro de 2026}
\rhead{\color{cinza}\footnotesize\bfseries Modelos com Efeitos Fixos e Interações de Mercado}
\lfoot{\color{cinza}\footnotesize\textit{Colinha do Apresentador --- Fernando Souza de Vieira}}
\rfoot{\color{verde}\bfseries Slide \slidenum\ de 15}
\renewcommand{\headrulewidth}{0.4pt}
\renewcommand{\footrulewidth}{0.4pt}

\tcbset{
  colback=cinzaclaro,
  colframe=verde,
  arc=2mm,
  boxrule=0.8pt,
  left=3.5mm,
  right=3.5mm,
  top=2.5mm,
  bottom=2.5mm
}

\begin{document}

\renewcommand{\slidenum}{1}
% ============================================================
% SLIDE 01 / 15
% ============================================================
\noindent
\begin{minipage}{\textwidth}
\begin{center}
{\color{cinza}\footnotesize\textbf{TÓPICOS EM MICROECONOMIA APLICADA --- 16 DE NOVEMBRO DE 2026}\\[1.2mm]}
{\color{cinza}\bfseries\Large SLIDE 01 / 15\\[1.5mm]}
{\color{verde}\bfseries\LARGE Abertura e Evolução da Pesquisa\\[2.5mm]}
\rule{\textwidth}{1pt}
\end{center}
\end{minipage}
\vspace{1.5mm}

\begin{tcolorbox}[title={\bfseries\small ÂNCORAS EMPÍRICAS E CONCEITUAIS DO SLIDE}]
\begin{itemize}[leftmargin=*,itemsep=1mm,topsep=0.5mm]
    \item \textbf{Contexto:} Disciplina \textit{Tópicos em Microeconomia Aplicada} (16/11/2026).
    \item \textbf{Evolução:} Continuidade direta da apresentação preliminar de 27/10/2026 (OLS).
    \item \textbf{Escopo:} 15 slides cobrindo intuição e hipóteses, lógica intuitiva do que foi testado e resultados consolidados do Capítulo 1 (Seções 3 e 4).
\end{itemize}
\end{tcolorbox}
\vspace{2mm}

\noindent{\color{verde}\large\bfseries Roteiro de Apresentação Oral}
\vspace{0.8mm}
\begin{enumerate}[leftmargin=*,label=\textbf{(\arabic*)},itemsep=2.8mm,topsep=1mm]
    \item Bom dia a todos. Hoje apresento a evolução da minha pesquisa sobre os determinantes econômicos e institucionais da conversão da vegetação nativa em Mato Grosso entre 2001 e 2024.
    \item Na apresentação anterior, de 27 de outubro, mostrei as regressões preliminares por Mínimos Quadrados Ordinários (OLS), discutindo seus resultados contra-intuitivos e suas limitações conceituais.
    \item Hoje trago o salto metodológico e empírico do Capítulo 1 consolidado: relembro primeiro a intuição, o problema e as hipóteses da pesquisa; explico intuitivamente o que foi testado e como os modelos de efeitos fixos resolvem os gargalos anteriores; e, por fim, interpreto os principais resultados sobre tecnologia, embargos ambientais, logística, regimes regulatórios e vazamento entre biomas.
\end{enumerate}

\newpage

\renewcommand{\slidenum}{2}
% ============================================================
% SLIDE 02 / 15
% ============================================================
\noindent
\begin{minipage}{\textwidth}
\begin{center}
{\color{cinza}\footnotesize\textbf{TÓPICOS EM MICROECONOMIA APLICADA --- 16 DE NOVEMBRO DE 2026}\\[1.2mm]}
{\color{cinza}\bfseries\Large SLIDE 02 / 15\\[1.5mm]}
{\color{verde}\bfseries\LARGE Relembrando a pesquisa: intuição, problema e hipóteses\\[2.5mm]}
\rule{\textwidth}{1pt}
\end{center}
\end{minipage}
\vspace{1.5mm}

\begin{tcolorbox}[title={\bfseries\small ÂNCORAS EMPÍRICAS E CONCEITUAIS DO SLIDE}]
\begin{itemize}[leftmargin=*,itemsep=1mm,topsep=0.5mm]
    \item \textbf{Intuição:} Decisão na fronteira opõe retorno de mercado ($\uparrow$ preços $\times$ solo fértil) à contenção institucional ($\downarrow$ IBAMA, embargos e acordos).
    \item \textbf{Mato Grosso (2001--2024):} Amazônia (Reserva Legal 80\%), Cerrado (35\%) e Pantanal (20\%) sob três marcos regulatórios.
    \item \textbf{H1:} Preços próprios aumentam conversão bruta, \textbf{amplificados pela aptidão} (\textbf{H1a}: $\beta > 0$) e \textbf{atenuados pela fiscalização} (\textbf{H1b}: $\delta < 0$).
\end{itemize}
\end{tcolorbox}
\vspace{2mm}

\noindent{\color{verde}\large\bfseries Roteiro de Apresentação Oral}
\vspace{0.8mm}
\begin{enumerate}[leftmargin=*,label=\textbf{(\arabic*)},itemsep=2.8mm,topsep=1mm]
    \item Antes de entrar nas novas regressões, vale relembrar rapidamente do que trata a pesquisa. Na fronteira agropecuária de Mato Grosso, o produtor decide entre manter a vegetação nativa ou convertê-la respondendo à tensão entre duas forças: o estímulo econômico de mercado e a contenção regulatória.
    \item Pelo lado econômico, quando o preço internacional de uma commodity sobe, o ganho potencial por hectare não cresce igual em todo lugar: ele cresce muito mais onde o solo é plano, profundo e mecanizável. Pelo lado institucional, a presença do IBAMA, os embargos e restrições comerciais aumentam o custo esperado de abrir irregularmente novas áreas.
    \item Mato Grosso é um laboratório ideal entre 2001 e 2024 porque combina liderança produtiva em soja, milho e pecuária com três biomas submetidos a regras de Reserva Legal muito diferentes: 80\% na Amazônia, 35\% no Cerrado e 20\% no Pantanal.
    \item Nossa pergunta central busca medir exatamente quanto os choques de preços induzem a conversão bruta em terras mais aptas e como a fiscalização condiciona esse efeito.
    \item A hipótese H1 desdobra-se em dois testes observáveis: H1a prevê amplificação pelos retornos potenciais do solo, e H1b prevê atenuação pela atuação ambiental. Tudo isso é confrontado entre lavouras versus pastagens e entre Amazônia versus Cerrado.
\end{enumerate}

\newpage

\renewcommand{\slidenum}{3}
% ============================================================
% SLIDE 03 / 15
% ============================================================
\noindent
\begin{minipage}{\textwidth}
\begin{center}
{\color{cinza}\footnotesize\textbf{TÓPICOS EM MICROECONOMIA APLICADA --- 16 DE NOVEMBRO DE 2026}\\[1.2mm]}
{\color{cinza}\bfseries\Large SLIDE 03 / 15\\[1.5mm]}
{\color{verde}\bfseries\LARGE Ponto de partida: o quebra-cabeça da OLS preliminar (27/10)\\[2.5mm]}
\rule{\textwidth}{1pt}
\end{center}
\end{minipage}
\vspace{1.5mm}

\begin{tcolorbox}[title={\bfseries\small ÂNCORAS EMPÍRICAS E CONCEITUAIS DO SLIDE}]
\begin{itemize}[leftmargin=*,itemsep=1mm,topsep=0.5mm]
    \item \textbf{OLS (27/10):} Preço soja M01 $= -0{,}0674\;(0{,}0667)$ pp; aptidão $= -0{,}0367\;(0{,}0222)$ pp ($R^2 \in [1{,}9\%;\, 5{,}6\%]$).
    \item \textbf{Por que falhou:} Viés de omissão espacial/temporal (áreas maduras sem estoque; pós-2005 juntou preços altos + PPCDAm), ausência de $P \times Z$ e 48\% zeros administrativos nos autos.
    \item \textbf{Passo adiante:} Limpar 612 imputações (18,1\% da base bruta) e identificar por interação dentro do painel.
\end{itemize}
\end{tcolorbox}
\vspace{2mm}

\noindent{\color{verde}\large\bfseries Roteiro de Apresentação Oral}
\vspace{0.8mm}
\begin{enumerate}[leftmargin=*,label=\textbf{(\arabic*)},itemsep=2.8mm,topsep=1mm]
    \item Agora, relembrando onde paramos na apresentação de 27 de outubro: quando rodamos modelos OLS simples e múltiplos, encontramos um aparente quebra-cabeça empírico. Tanto o preço da soja quanto a aptidão agronômica apareciam com sinal negativo na conversão para lavouras, sem significância estatística e com R quadrado entre 2\% e 6\%.
    \item Vimos que isso não refutava a teoria, mas expunha três gargalos da modelagem preliminar. Primeiro, um grave viés de variável omitida no espaço e no tempo: por um lado, municípios muito férteis já desmataram no passado e têm menos estoque restante; por outro, no tempo, a disparada de preços após 2005 coincidiu com políticas fortes de combate ao desmatamento (PPCDAm e Moratória), criando correlação negativa espúria na OLS empilhada.
    \item Segundo, preços e aptidão entravam apenas somados, quando a teoria econômica prevê que o preço só altera a atratividade da conversão de forma multiplicativa pela produtividade local da terra.
    \item Terceiro, ao incluirmos autos do IBAMA, perdíamos mais da metade da amostra observada, e na base original havia imputações e zeros administrativos que distorciam as estimativas iniciais. Precisávamos mudar o desenho empírico.
\end{enumerate}

\newpage

\renewcommand{\slidenum}{4}
% ============================================================
% SLIDE 04 / 15
% ============================================================
\noindent
\begin{minipage}{\textwidth}
\begin{center}
{\color{cinza}\footnotesize\textbf{TÓPICOS EM MICROECONOMIA APLICADA --- 16 DE NOVEMBRO DE 2026}\\[1.2mm]}
{\color{cinza}\bfseries\Large SLIDE 04 / 15\\[1.5mm]}
{\color{verde}\bfseries\LARGE O que foi testado e como: intuição das regressões em painel\\[2.5mm]}
\rule{\textwidth}{1pt}
\end{center}
\end{minipage}
\vspace{1.5mm}

\begin{tcolorbox}[title={\bfseries\small ÂNCORAS EMPÍRICAS E CONCEITUAIS DO SLIDE}]
\begin{itemize}[leftmargin=*,itemsep=1mm,topsep=0.5mm]
    \item \textbf{Resumo em 1 frase:} Efeitos fixos (1 e 2) nivelam diferenças permanentes entre municípios ($\alpha_m$) e a maré anual do bioma ($\lambda_{b(m),t}$) para que as interações (3 e 4) isolem o \textbf{acelerador de mercado} ($P \times Z$) e o \textbf{freio ambiental} ($P \times Z \times \text{Inst}$).
    \item \textbf{5 baterias (178 modelos/testes):} (1) Núcleo H1a; (2) Pacote técnico/2ª safra; (3) Logística e embargos H1b; (4) Regimes e vazamento; (5) Pastagem e placebo $t+1$.
\end{itemize}
\end{tcolorbox}
\vspace{2mm}

\noindent{\color{verde}\large\bfseries Roteiro de Apresentação e Explicação Passo a Passo (Nível Graduação)}
\vspace{0.8mm}
{\small
\begin{enumerate}[leftmargin=*,label=\textbf{(\arabic*)},itemsep=1.2mm,topsep=1mm]
    \item \textbf{Introdução:} Antes das equações formais, este slide mostra intuitivamente o que testamos e \textit{como} comparamos os municípios sem cair nas armadilhas da regressão simples (OLS).
    \item \textbf{Por que uma regressão comum engana?} Na OLS em ``foto'', duas armadilhas quebram o \textit{ceteris paribus}: \textit{(i) no espaço}, cidades férteis e maduras (como Sorriso) já foram muito abertas no passado e têm pouca mata restante, enquanto fronteiras novas ainda têm floresta --- comparar uma com a outra faz parecer que solo bom desmata \textit{menos}; \textit{(ii) no tempo}, após 2005 o preço da soja subiu junto com o aperto do PPCDAm e da Moratória da Soja, gerando associação espúria negativa.
    \item \textbf{Passo 1 --- Município contra ele mesmo ($\alpha_m$):} Em vez de comparar Sorriso com uma cidade do Pantanal, damos a cada município seu próprio ponto de partida e comparamos o \textit{filme} daquela cidade contra sua própria média histórica (2001--2024). Tudo o que é fixo (relevo, distância, história fundiária) é descontado: \textit{``Quando o preço muda, o município X desmata mais ou menos que o normal dele mesmo?''}
    \item \textbf{Passo 2 --- Mesmo bioma no mesmo ano ($\lambda_{b(m),t}$):} Descontamos a ``maré macroeconômica do ano'' em cada bioma (dólar, Selic, El Niño, leis federais). Se o Cerrado inteiro desacelerou 20\% num ano de crise, subtraímos esses 20\% de todos os municípios do Cerrado naquele ano. O tabuleiro fica nivelado para ver quem fugiu da média em relação aos vizinhos.
    \item \textbf{Passo 3 --- O teste de H1a ($P \times Z$, o acelerador):} Como o preço internacional sobe igual para todo o estado no mesmo ano, multiplicamos preço por aptidão ($Z$), pois a receita extra por hectare é $\Delta\text{Preço} \times \text{Produtividade}$. Pergunta: \textit{``Quando a soja sobe lá fora, o município de solo fértil acelera a derrubada mais que o vizinho de solo fraco do mesmo bioma?''} (Tabela 1: sim, $+0{,}0848^{**}$).
    \item \textbf{Passo 4 --- O teste de H1b ($P \times Z \times \text{Inst}$, o freio):} Diante da mesma alta externa entre duas cidades igualmente férteis do mesmo bioma, uma sofreu embargos do IBAMA no ano anterior (travando crédito e venda) e a outra não. Pergunta: \textit{``Onde houve embargo, o impulso de abrir solo fértil perdeu força?''} Se $P \times Z$ é o acelerador, essa interação mede o freio institucional (Tabela 3: $-0{,}0314^{*}$, atenuando em 23,3\% no 1º embargo).
    \item \textbf{Em resumo e mapa do que foi testado:} Os passos 1 e 2 limpam diferenças permanentes entre cidades e a maré de cada ano para que 3 e 4 isolem o acelerador de mercado e o freio ambiental. Do lado direito está o roteiro das 5 baterias (Tabelas 1 a 5): núcleo H1a; pacote técnico e 2ª safra; logística e embargos vs.\ autos; regimes e Amazônia vs.\ Cerrado; pastagem e placebo $t+1$.
\end{enumerate}
}

\newpage

\renewcommand{\slidenum}{5}
% ============================================================
% SLIDE 05 / 15
% ============================================================
\noindent
\begin{minipage}{\textwidth}
\begin{center}
{\color{cinza}\footnotesize\textbf{TÓPICOS EM MICROECONOMIA APLICADA --- 16 DE NOVEMBRO DE 2026}\\[1.2mm]}
{\color{cinza}\bfseries\Large SLIDE 05 / 15\\[1.5mm]}
{\color{verde}\bfseries\LARGE Especificação econométrica formal: efeitos fixos e interações\\[2.5mm]}
\rule{\textwidth}{1pt}
\end{center}
\end{minipage}
\vspace{1.5mm}

\begin{tcolorbox}[title={\bfseries\small ÂNCORAS EMPÍRICAS E CONCEITUAIS DO SLIDE}]
\begin{itemize}[leftmargin=*,itemsep=1mm,topsep=0.5mm]
    \item \textbf{Eq.\ (H1a):} $Y_{mt}^d = \alpha_m + \lambda_{b(m),t} + \beta_d(\ln P_{jt-1} \times Z_{jm}) + \gamma_d(Z_{jm} \times \tau_t) + \rho_d R_{mt} + \varepsilon_{mt}^d$.
    \item \textbf{Eq.\ (H1b):} Acrescenta $\delta_d(\ln P_{jt-1} \times Z_{jm} \times \text{Inst}_{mt-1})$ e termos inferiores.
    \item \textbf{Controles-chave:} Tendência $Z \times \tau_t$ blinda $\beta_d$ contra difusão tecnológica secular; $R_{mt}$ (CHIRPS); cluster por município ($G=134$).
\end{itemize}
\end{tcolorbox}
\vspace{2mm}

\noindent{\color{verde}\large\bfseries Roteiro de Apresentação Oral}
\vspace{0.8mm}
\begin{enumerate}[leftmargin=*,label=\textbf{(\arabic*)},itemsep=2.8mm,topsep=1mm]
    \item Traduzindo essa intuição para a notação formal do artigo: estas são as duas equações centrais da Seção 3, estimadas por painel com efeitos fixos bidirecionais.
    \item Na equação de H1a, os efeitos fixos municipais (alfa) absorvem o nível fixo da aptidão Z, o relevo e toda característica geográfica permanente; os efeitos fixos de bioma por ano (lambda) absorvem o nível comum do preço internacional e os choques agregados de cada ano por bioma.
    \item Por isso, a identificação vem inteiramente da interação entre o log do preço defasado e a aptidão padronizada Z, esperando-se beta positivo sob H1a. Também controlamos explicitamente pela tendência linear diferenciada por aptidão (Z vezes tau) para evitar que beta capture apenas a modernização tecnológica gradual do Centro-Oeste, e pela anomalia de chuva contemporânea (R).
    \item Na extensão H1b, acrescentamos a interação tripla com a variável institucional defasada (autos de infração ou embargos do IBAMA) junto de todos os termos de ordem inferior não absorvidos, esperando delta negativo quando a fiscalização freia a resposta aos preços. A inferência base agrupa os erros-padrão pelos 134 municípios efetivos.
\end{enumerate}

\newpage

\renewcommand{\slidenum}{6}
% ============================================================
% SLIDE 06 / 15
% ============================================================
\noindent
\begin{minipage}{\textwidth}
\begin{center}
{\color{cinza}\footnotesize\textbf{TÓPICOS EM MICROECONOMIA APLICADA --- 16 DE NOVEMBRO DE 2026}\\[1.2mm]}
{\color{cinza}\bfseries\Large SLIDE 06 / 15\\[1.5mm]}
{\color{verde}\bfseries\LARGE Construção da base: painel desbalanceado (2001--2024)\\[2.5mm]}
\rule{\textwidth}{1pt}
\end{center}
\end{minipage}
\vspace{1.5mm}

\begin{tcolorbox}[title={\bfseries\small ÂNCORAS EMPÍRICAS E CONCEITUAIS DO SLIDE}]
\begin{itemize}[leftmargin=*,itemsep=1mm,topsep=0.5mm]
    \item \textbf{Painel efetivo:} $\text{3.384 (teórico)} - 612\text{ (imputados)} - 5\text{ (singletons)} = \mathbf{2.767\text{ obs.}}$ em \textbf{134 municípios}.
    \item \textbf{Subamostra fiscal (autos):} $1.135$ obs.\ em $123$ municípios; já os embargos geocodificados ($11.513$ termos) preservam toda a amostra de $2.767$.
    \item \textbf{Média de $Y^{\text{agr}}$:} $0{,}112$ pp/ano (conversão bruta MapBiomas 10.1); $1\text{ DP}$ de $Z \approx 380$ kg/ha (soja) e $779$ kg/ha (milho). Preço da terra e crédito excluídos (\textit{bad controls}).
\end{itemize}
\end{tcolorbox}
\vspace{2mm}

\noindent{\color{verde}\large\bfseries Roteiro de Apresentação Oral}
\vspace{0.8mm}
\begin{enumerate}[leftmargin=*,label=\textbf{(\arabic*)},itemsep=2.8mm,topsep=1mm]
    \item Aqui detalho a base que alimenta todas as regressões de efeitos fixos. A grade completa de Mato Grosso possui 141 municípios em 24 anos, somando 3.384 pares teóricos.
    \item Na versão consolidada, eliminamos integralmente as 612 observações imputadas que existiam na base precária, chegando a 2.772 pares observados; ao estimar com efeitos fixos de município e bioma por ano, o fixest descarta 5 observações unitárias, resultando em exatamente 2.767 observações efetivas em 134 municípios.
    \item Já os embargos ambientais geocodificados (mais de 11 mil registros com coordenadas válidas) permitem estimar a modulação institucional nessa amostra plena de 2.767 observações, ao contrário dos autos administrativos de flora, cuja subamostra elegível restringe-se a 1.135 observações em 123 municípios efetivos.
    \item Nossa variável dependente Y é a conversão bruta anual do MapBiomas 10.1 — sem compensação contábil por áreas em regeneração — com média observada de 0,112 ponto percentual ao ano. E mantivemos fora da regressão o preço de mercado da terra e o crédito contemporâneo por serem maus controles (bad controls) determinados em equilíbrio pelo próprio choque de commodities.
\end{enumerate}

\newpage

\renewcommand{\slidenum}{7}
% ============================================================
% SLIDE 07 / 15
% ============================================================
\noindent
\begin{minipage}{\textwidth}
\begin{center}
{\color{cinza}\footnotesize\textbf{TÓPICOS EM MICROECONOMIA APLICADA --- 16 DE NOVEMBRO DE 2026}\\[1.2mm]}
{\color{cinza}\bfseries\Large SLIDE 07 / 15\\[1.5mm]}
{\color{verde}\bfseries\LARGE Resultados principais: confirmação de H1a na agricultura\\[2.5mm]}
\rule{\textwidth}{1pt}
\end{center}
\end{minipage}
\vspace{1.5mm}

\begin{tcolorbox}[title={\bfseries\small ÂNCORAS EMPÍRICAS E CONCEITUAIS DO SLIDE}]
\begin{itemize}[leftmargin=*,itemsep=1mm,topsep=0.5mm]
    \item \textbf{Tabela 1 (H1a):} Soja M1 $= 0{,}0848^{**}\;(0{,}0358)$, $p=0{,}019$; Milho M2 $= 0{,}0519^{***}\;(0{,}0179)$, $p=0{,}004$.
    \item \textbf{Estimação conjunta M3:} Preços altamente correlacionados ($r = 0{,}948$), mas Wald conjunto rejeita nulidade ($F = 4{,}12$, $p = 0{,}018$).
    \item \textbf{Magnitude:} $+10\%$ no preço da soja $\Rightarrow +0{,}0081$ pp (+7,2\%); $+50\% \Rightarrow +0{,}0344$ pp (+30,6\% sobre a média, ou $\approx 125$ ha/ano).
\end{itemize}
\end{tcolorbox}
\vspace{2mm}

\noindent{\color{verde}\large\bfseries Roteiro de Apresentação Oral}
\vspace{0.8mm}
\begin{enumerate}[leftmargin=*,label=\textbf{(\arabic*)},itemsep=2.8mm,topsep=1mm]
    \item Este slide traz o resultado empírico central da tese, correspondente à Tabela 1 do artigo.
    \item Vejam a inversão completa em relação ao OLS de 27 de outubro: ao controlarmos os atributos fixos municipais e os choques de bioma por ano, o coeficiente da interação de preços torna-se positivo e estatisticamente significante.
    \item Na coluna 1, para a soja, a estimativa é de 0,0848, significativa a 5\%. Na coluna 2, para o milho, é de 0,0519, significativa a 1\%. Reparem também no papel da tendência Z vezes tau, positiva e significante: sem ela, beta subiria artificialmente para 0,142 na soja por absorver a modernização gradual do estado.
    \item Na coluna 3, ao incluirmos soja e milho juntos, a altíssima correlação entre seus preços internacionais (r = 0,948) divide a precisão individual, mas o teste de Wald conjunto rejeita enfaticamente que ambos os efeitos sejam nulos (F = 4,12, p = 0,018).
    \item Em termos práticos: uma alta de 50\% no preço real da soja gera um diferencial de 0,0344 ponto percentual na conversão anual entre municípios separados por 1 desvio-padrão de aptidão — equivalente a mais de 30\% da taxa média estadual, ou cerca de 125 hectares adicionais por ano em um estoque médio. H1a é claramente confirmada na associação condicional.
\end{enumerate}

\newpage

\renewcommand{\slidenum}{8}
% ============================================================
% SLIDE 08 / 15
% ============================================================
\noindent
\begin{minipage}{\textwidth}
\begin{center}
{\color{cinza}\footnotesize\textbf{TÓPICOS EM MICROECONOMIA APLICADA --- 16 DE NOVEMBRO DE 2026}\\[1.2mm]}
{\color{cinza}\bfseries\Large SLIDE 08 / 15\\[1.5mm]}
{\color{verde}\bfseries\LARGE O mecanismo tecnológico: a aptidão depende do pacote técnico\\[2.5mm]}
\rule{\textwidth}{1pt}
\end{center}
\end{minipage}
\vspace{1.5mm}

\begin{tcolorbox}[title={\bfseries\small ÂNCORAS EMPÍRICAS E CONCEITUAIS DO SLIDE}]
\begin{itemize}[leftmargin=*,itemsep=1mm,topsep=0.5mm]
    \item \textbf{Tabela 2:} Baixo insumo ($-0{,}0191$, ns), Adequação climática ($-0{,}0874^{**}$) e PAM 91--00 ($-0{,}0644$, ns; $N=1.016$, $45$ mun.) não geram gradiente positivo.
    \item \textbf{Ocupação pré-2001 ($H_i$):} Ortogonal ao GAEZ ($r = -0{,}025$); associa-se a \textbf{menor} conversão bruta ($-0{,}1160^{***}$).
    \item \textbf{Intensificação (col.\ 7):} Municípios já consolidados expandiram milho de 2ª safra em $+2{,}51^{***}$ pp/década ($R^2$ \textit{within} $= 0{,}324$), poupando mata.
\end{itemize}
\end{tcolorbox}
\vspace{2mm}

\noindent{\color{verde}\large\bfseries Roteiro de Apresentação Oral}
\vspace{0.8mm}
\begin{enumerate}[leftmargin=*,label=\textbf{(\arabic*)},itemsep=2.8mm,topsep=1mm]
    \item Na Tabela 2, perguntamos: qualquer medida de solo gera esse efeito positivo ou ele depende especificamente da agricultura empresarial de alta tecnologia?
    \item As colunas 1 a 3 dão uma resposta incisiva: quando substituímos a aptidão de alto insumo pelo cenário GAEZ de baixo insumo, pelo índice genérico de adequação agroclimática ou pela produtividade comercial histórica da PAM nos anos 90, o gradiente positivo desaparece completamente ou até inverte.
    \item Isso faz todo sentido agronômico: solos de Cerrado em estado natural são ácidos e pouco férteis. O choque de preços internacionais só induz abertura extensiva de lavouras onde é rentável aplicar correção pesada de solo, fertilizantes e mecanização.
    \item Além disso, as colunas 4 a 7 revelam a dualidade territorial de Mato Grosso: a ocupação agrícola histórica (1985–2000) é ortogonal ao potencial GAEZ (r = -0,025) e apresenta coeficiente fortemente negativo sobre a conversão bruta (-0,1160). Onde a agricultura já estava consolidada, os produtores responderam aos altos preços intensificando a área plantada com a segunda safra de milho (+2,51 pontos percentuais por década na coluna 7), poupando vegetação nativa enquanto a abertura bruta deslocava-se para as novas fronteiras.
\end{enumerate}

\newpage

\renewcommand{\slidenum}{9}
% ============================================================
% SLIDE 09 / 15
% ============================================================
\noindent
\begin{minipage}{\textwidth}
\begin{center}
{\color{cinza}\footnotesize\textbf{TÓPICOS EM MICROECONOMIA APLICADA --- 16 DE NOVEMBRO DE 2026}\\[1.2mm]}
{\color{cinza}\bfseries\Large SLIDE 09 / 15\\[1.5mm]}
{\color{verde}\bfseries\LARGE Dissuasão ambiental (H1b): o poder dos embargos do IBAMA\\[2.5mm]}
\rule{\textwidth}{1pt}
\end{center}
\end{minipage}
\vspace{1.5mm}

\begin{tcolorbox}[title={\bfseries\small ÂNCORAS EMPÍRICAS E CONCEITUAIS DO SLIDE}]
\begin{itemize}[leftmargin=*,itemsep=1mm,topsep=0.5mm]
    \item \textbf{Autos vs.\ Embargos (Tabela 3):} Autos de flora ($N=1.135$, $123$ mun.) são nulos ($0{,}0002$; tripla em M7 $p=0{,}264$) por morosidade pecuniária e menor amostra.
    \item \textbf{Área embargada ($t-1$):} $-0{,}0076^{***}\;(0{,}0029)$ reduz diretamente a conversão agrícola ($N=2.767$).
    \item \textbf{Modulação H1b (col.\ 6):} Tripla $-0{,}0314^{*}\;(0{,}0179)$, $p=0{,}081$. 1º embargo ($\ln 2$) reduz o gradiente em \textbf{23,3\%} ($-33{,}6\%$ por log) via bloqueio de terra e crédito.
\end{itemize}
\end{tcolorbox}
\vspace{2mm}

\noindent{\color{verde}\large\bfseries Roteiro de Apresentação Oral}
\vspace{0.8mm}
\begin{enumerate}[leftmargin=*,label=\textbf{(\arabic*)},itemsep=2.8mm,topsep=1mm]
    \item Na Tabela 3, colunas 4 a 7, testamos diretamente a hipótese H1b: instrumentos de fiscalização ambiental conseguem moderar a sensibilidade da conversão aos choques de preços?
    \item O contraste entre os instrumentos é um dos achados mais importantes da pesquisa. Na coluna 7, utilizando os autos de infração de flora do IBAMA na subamostra fiscal elegível (1.135 observações em 123 municípios efetivos), o coeficiente é estatisticamente nulo (0,0002; e também nulo quando testado em tripla interação). Como discutimos no texto, isso reflete tanto a baixa liquidez e lentidão judicial das multas pecuniárias quanto a redução amostral e de variação nessa subamostra.
    \item Em contrapartida, ao utilizarmos os embargos ambientais geocodificados — que preservam toda a amostra observada de 2.767 observações —, a área embargada defasada reduz diretamente a conversão a 1\% de significância (-0,0076).
    \item Ainda mais decisivo para H1b: na coluna 6, a interação tripla entre preço, aptidão e número de novos embargos entra negativa (-0,0314, significativa a 10\%). O registro do primeiro embargo em um município está associado a uma atenuação estimada de 23,3\% no gradiente preço-aptidão (e uma taxa marginal de -33,6\% por unidade logarítmica).
    \item A explicação econômica é clara: diferentemente da multa avulsa, o embargo ambiental impõe restrição patrimonial imediata à área autuada e bloqueia o acesso ao crédito rural oficial, elevando o custo esperado de responder a altas de preços via desmatamento.
\end{enumerate}

\newpage

\renewcommand{\slidenum}{10}
% ============================================================
% SLIDE 10 / 15
% ============================================================
\noindent
\begin{minipage}{\textwidth}
\begin{center}
{\color{cinza}\footnotesize\textbf{TÓPICOS EM MICROECONOMIA APLICADA --- 16 DE NOVEMBRO DE 2026}\\[1.2mm]}
{\color{cinza}\bfseries\Large SLIDE 10 / 15\\[1.5mm]}
{\color{verde}\bfseries\LARGE Logística e infraestrutura de transporte (Tabela 3)\\[2.5mm]}
\rule{\textwidth}{1pt}
\end{center}
\end{minipage}
\vspace{1.5mm}

\begin{tcolorbox}[title={\bfseries\small ÂNCORAS EMPÍRICAS E CONCEITUAIS DO SLIDE}]
\begin{itemize}[leftmargin=*,itemsep=1mm,topsep=0.5mm]
    \item \textbf{Núcleo preservado:} Com logística, soja $= 0{,}0940^{**}\;(0{,}0394)$ e milho $= 0{,}0656^{***}\;(0{,}0214)$.
    \item \textbf{Distância mínima a portos ($D_{it}^{\text{min}}$):} $+0{,}0252^{***}\;(0{,}0079)$ na soja e $+0{,}0261^{***}\;(0{,}0084)$ no milho (após Arco Norte em 2014).
    \item \textbf{Ferrovia Rondonópolis 2013 ($F_{it}^{2013}$):} $-0{,}1689^{***}\;(0{,}0453)$ na soja; barateou frete em área madura, estimulando intensificação em vez de abertura pioneira.
\end{itemize}
\end{tcolorbox}
\vspace{2mm}

\noindent{\color{verde}\large\bfseries Roteiro de Apresentação Oral}
\vspace{0.8mm}
\begin{enumerate}[leftmargin=*,label=\textbf{(\arabic*)},itemsep=2.8mm,topsep=1mm]
    \item Ainda na Tabela 3, nas colunas 1 a 3, avaliamos o papel da infraestrutura logística — calculada diretamente por distâncias geodésicas dos centroides municipais aos portos do Sul, ao Arco Norte após 2014 e ao terminal ferroviário de Rondonópolis, inaugurado em 2013.
    \item Primeiro ponto: incluir logística não enfraquece nosso achado central de H1a; pelo contrário, soja (0,0940) e milho (0,0656) ficam até mais precisos e fortes.
    \item Segundo ponto: a inauguração da ferrovia em Rondonópolis em 2013 exibe forte associação negativa com a conversão bruta (-0,1689 na soja e -0,1280 no milho). Ao mesmo tempo, a distância mínima aos portos tem coeficiente positivo (0,0252) depois que o Arco Norte encurtou as rotas do norte mato-grossense após 2014.
    \item Qual é a intuição microeconômica desse resultado? Seguindo von Thünen, o produtor reage ao preço internacional líquido do frete local. Como a ferrovia chegou a uma região agrícola já consolidada no sul-sudeste do estado, o barateamento do frete aumentou o custo de oportunidade de abrir mata virgem distante, incentivando a intensificação nas lavouras já estabelecidas em detrimento da fronteira pioneira.
\end{enumerate}

\newpage

\renewcommand{\slidenum}{11}
% ============================================================
% SLIDE 11 / 15
% ============================================================
\noindent
\begin{minipage}{\textwidth}
\begin{center}
{\color{cinza}\footnotesize\textbf{TÓPICOS EM MICROECONOMIA APLICADA --- 16 DE NOVEMBRO DE 2026}\\[1.2mm]}
{\color{cinza}\bfseries\Large SLIDE 11 / 15\\[1.5mm]}
{\color{verde}\bfseries\LARGE Regimes regulatórios: quebra estrutural pós-2005\\[2.5mm]}
\rule{\textwidth}{1pt}
\end{center}
\end{minipage}
\vspace{1.5mm}

\begin{tcolorbox}[title={\bfseries\small ÂNCORAS EMPÍRICAS E CONCEITUAIS DO SLIDE}]
\begin{itemize}[leftmargin=*,itemsep=1mm,topsep=0.5mm]
    \item \textbf{Tabela 4 — Painel A (MT):} Boom 2001--05 ($\beta_s = 0{,}3541^{***}$, $>4\times$ a média histórica; $\beta_m = 0{,}4234^{***}$); pós-2005 cai a zero ($0{,}0108$ e $0{,}0120$).
    \item \textbf{Marcos regulatórios:} PPCDAm/satélites (2004), Moratória da Soja (2006), crédito restrito e Código Florestal/CAR (2012).
    \item \textbf{Ressalva metodológica (Painel B):} A magnitude de 2001--05 varia conforme parametrização de tendências/níveis, mas o desacoplamento médio pós-2005 mantém-se firme.
\end{itemize}
\end{tcolorbox}
\vspace{2mm}

\noindent{\color{verde}\large\bfseries Roteiro de Apresentação Oral}
\vspace{0.8mm}
\begin{enumerate}[leftmargin=*,label=\textbf{(\arabic*)},itemsep=2.8mm,topsep=1mm]
    \item Na Tabela 4, testamos se a relação entre preços internacionais e conversão manteve a mesma força ao longo dos 24 anos ou se houve quebra estrutural com a mudança das regras ambientais.
    \item Dividimos o painel em três regimes institucionais: 2001–2005 (abertura do superciclo e desmatamento recorde), 2006–2011 (endurecimento do PPCDAm, Moratória da Soja e restrição ao crédito) e 2012–2024 (Novo Código Florestal e governança cadastral pelo CAR).
    \item No Painel A, os resultados revelam uma ruptura nítida: toda a sensibilidade positiva concentra-se no quinquênio 2001–2005, com coeficientes superiores a 0,35 na soja e 0,42 no milho (mais de quatro vezes a média de 24 anos). A partir de 2006, o coeficiente estadual cai para perto de zero e perde significância estatística.
    \item Com rigor metodológico, também reportamos no Painel B (slide seguinte) uma especificação irrestrita permitindo níveis diferentes por regime: nela, a magnitude média inicial de 2001–2005 é sensível à parametrização das tendências, mas o desacoplamento médio do pós-2005 na fronteira agrícola se mantém firme.
    \item Resta perguntar: esse desacoplamento após 2005 aconteceu igualmente na Amazônia e no Cerrado?
\end{enumerate}

\newpage

\renewcommand{\slidenum}{12}
% ============================================================
% SLIDE 12 / 15
% ============================================================
\noindent
\begin{minipage}{\textwidth}
\begin{center}
{\color{cinza}\footnotesize\textbf{TÓPICOS EM MICROECONOMIA APLICADA --- 16 DE NOVEMBRO DE 2026}\\[1.2mm]}
{\color{cinza}\bfseries\Large SLIDE 12 / 15\\[1.5mm]}
{\color{verde}\bfseries\LARGE Amazônia vs. Cerrado: evidência de vazamento regulatório\\[2.5mm]}
\rule{\textwidth}{1pt}
\end{center}
\end{minipage}
\vspace{1.5mm}

\begin{tcolorbox}[title={\bfseries\small ÂNCORAS EMPÍRICAS E CONCEITUAIS DO SLIDE}]
\begin{itemize}[leftmargin=*,itemsep=1mm,topsep=0.5mm]
    \item \textbf{Tabela 4 — Painel B (2006--2011):} Amazônia ($1.303$ obs., $71$ mun.) $= -0{,}0704^{**}\;(0{,}0334)$; Cerrado ($1.412$ obs., $60$ mun.) $= +0{,}0569^{**}\;(0{,}0258)$.
    \item \textbf{Assimetria institucional:} RL 80\% + Moratória da Soja na Amazônia vs.\ RL 35\% sem Moratória no Cerrado (estável entre Painéis A e B).
    \item \textbf{Leitura científica:} Evidência compatível com desacoplamento amazônico e vazamento (\textit{leakage}) de pressão para o Cerrado vizinho.
\end{itemize}
\end{tcolorbox}
\vspace{2mm}

\noindent{\color{verde}\large\bfseries Roteiro de Apresentação Oral}
\vspace{0.8mm}
\begin{enumerate}[leftmargin=*,label=\textbf{(\arabic*)},itemsep=2.8mm,topsep=1mm]
    \item Este slide mostra um dos resultados mais expressivos do artigo, separando as subamostras de bioma predominante: 71 municípios amazônicos e 60 de Cerrado no Painel B (e idêntico em sinal e significância no Painel A).
    \item Lembram que a Moratória da Soja de 2006 proibiu a compra de grãos oriundos de desmatamentos novos apenas no bioma Amazônia? Olhem o que acontece exatamente na janela 2006–2011:
    \item Na Amazônia, a sensibilidade da conversão aos preços da soja inverte para menos 0,0704, significativo a 5\% — sinal de que o bloqueio de cadeia, somado ao PPCDAm e à Reserva Legal de 80\%, conteve o avanço direto da soja sobre a floresta mesmo em anos de alta externa.
    \item Já no Cerrado, no exato mesmo período de 2006 a 2011, o coeficiente manteve-se positivo em mais 0,0569, também significativo a 5\%.
    \item Embora identificar causalmente a transferência de cada produtor exija um desenho contrafactual próprio de spillover, essa divergência estável entre os dois painéis é plenamente compatível com a hipótese de vazamento regulatório (leakage): proteger unilateralmente um bioma sem estender regras semelhantes ao bioma vizinho redirecionou a expansão extensiva para as savanas do Cerrado.
\end{enumerate}

\newpage

\renewcommand{\slidenum}{13}
% ============================================================
% SLIDE 13 / 15
% ============================================================
\noindent
\begin{minipage}{\textwidth}
\begin{center}
{\color{cinza}\footnotesize\textbf{TÓPICOS EM MICROECONOMIA APLICADA --- 16 DE NOVEMBRO DE 2026}\\[1.2mm]}
{\color{cinza}\bfseries\Large SLIDE 13 / 15\\[1.5mm]}
{\color{verde}\bfseries\LARGE Destinos comparados: a pecuária como ponta de lança\\[2.5mm]}
\rule{\textwidth}{1pt}
\end{center}
\end{minipage}
\vspace{1.5mm}

\begin{tcolorbox}[title={\bfseries\small ÂNCORAS EMPÍRICAS E CONCEITUAIS DO SLIDE}]
\begin{itemize}[leftmargin=*,itemsep=1mm,topsep=0.5mm]
    \item \textbf{Tabela 5, col.\ 1:} Agricultura ($\beta_s = 0{,}0848^{**}$, $\beta_m = 0{,}0519^{***}$) vs.\ Pastagem ($\beta_s = 0{,}3206^{***}$, $\beta_m = 0{,}4600^{***}$).
    \item \textbf{Proporção:} Pastagem tem gradiente $3{,}8\times$ maior na soja e $8{,}9\times$ maior no milho.
    \item \textbf{Interpretação:} Pecuária como vetor fundiário preparatório de baixo custo (>90\% da abertura primária), além de refletir queda pós-2005 em áreas menos mecanizáveis.
\end{itemize}
\end{tcolorbox}
\vspace{2mm}

\noindent{\color{verde}\large\bfseries Roteiro de Apresentação Oral}
\vspace{0.8mm}
\begin{enumerate}[leftmargin=*,label=\textbf{(\arabic*)},itemsep=2.8mm,topsep=1mm]
    \item Na Tabela 5, coluna 1, testamos a mesma especificação tendo como desfecho a conversão bruta para pastagem, já que a análise descritiva mostrou que mais de 90\% da vegetação nativa suprimida em Mato Grosso vira pasto em um primeiro momento.
    \item Vejam o contraste de magnitude: o gradiente associado ao preço e à aptidão GAEZ na pastagem é de 0,3206 na soja e 0,4600 no milho — quase quatro vezes maior que o da lavoura direta na soja, e quase nove vezes maior no milho.
    \item Por que o preço internacional de soja e milho e a aptidão mecanizável para grãos se associam tão fortemente à abertura de pastagens?
    \item A literatura de fronteira agrícola e nossos resultados apontam para a dinâmica sequencial do uso da terra: abrir floresta diretamente para soja mecanizada exige destoca profunda, nivelamento e calagem caríssimos. A pecuária entra como ponta de lança de baixo custo para limpar o terreno, assegurar posse fundiária e esperar o assentamento do solo, sendo convertida anos depois em áreas de grãos.
    \item Com transparência analítica, também destacamos que esse forte sinal na pastagem sob a tendência linear 2001–2024 reflete parcialmente a retração generalizada do desmatamento pós-2005 justamente nas regiões menos aptas à mecanização, como boa parte da Amazônia mato-grossense.
\end{enumerate}

\newpage

\renewcommand{\slidenum}{14}
% ============================================================
% SLIDE 14 / 15
% ============================================================
\noindent
\begin{minipage}{\textwidth}
\begin{center}
{\color{cinza}\footnotesize\textbf{TÓPICOS EM MICROECONOMIA APLICADA --- 16 DE NOVEMBRO DE 2026}\\[1.2mm]}
{\color{cinza}\bfseries\Large SLIDE 14 / 15\\[1.5mm]}
{\color{verde}\bfseries\LARGE Rigor metodológico: sensibilidade e teste de placebo\\[2.5mm]}
\rule{\textwidth}{1pt}
\end{center}
\end{minipage}
\vspace{1.5mm}

\begin{tcolorbox}[title={\bfseries\small ÂNCORAS EMPÍRICAS E CONCEITUAIS DO SLIDE}]
\begin{itemize}[leftmargin=*,itemsep=1mm,topsep=0.5mm]
    \item \textbf{Robustez:} Sobrevive a cluster duplo (mun.+ano), winsorização 99\%, sem Pantanal, FWL-HAC e \textit{leave-one-year-out} (24/24 anos $p<0{,}10$); zera em 2006--2024 ($-0{,}0067$).
    \item \textbf{Placebo temporal ($t+1$):} Soja $= 0{,}0844^{**}\;(0{,}0338)$ e Milho $= 0{,}1026^{***}\;(0{,}0304)$ ($N=2.633$).
    \item \textbf{Interpretação honesta:} Autocorrelação de preços ($r \approx 0{,}73$ na soja), quebra de nível pós-2006 e preparo plurianual (2--3 anos) $\Rightarrow$ associações aos ciclos, não choques pontuais ano a ano.
\end{itemize}
\end{tcolorbox}
\vspace{2mm}

\noindent{\color{verde}\large\bfseries Roteiro de Apresentação Oral}
\vspace{0.8mm}
\begin{enumerate}[leftmargin=*,label=\textbf{(\arabic*)},itemsep=2.8mm,topsep=1mm]
    \item Na Tabela 5, colunas 2 a 7, submetemos nossos resultados a uma bateria rigorosa de testes de sensibilidade e limites de identificação.
    \item Do lado da robustez estatística, as estimativas sobrevivem à dupla clusterização por município e ano, ao truncamento de extremos por winsorização a 99\%, à transformação arco-seno hiperbólico, à regressão nos 24 pontos anuais (FWL-HAC) e à exclusão iterativa de cada ano da amostra (significativas em todos os 24 cortes). Também confirmam o desacoplamento na subjanela 2006–2024, onde o coeficiente zera.
    \item Contudo, a coluna 4 traz um teste fundamental de transparência metodológica: ao inserirmos o preço futuro em t+1 como teste de placebo, ele entra positivo e significativo (0,0844 na soja e 0,1026 no milho).
    \item Em vez de omitir esse resultado, nós o incorporamos diretamente à interpretação do paper: primeiro, as séries de preços de commodities exibem forte autocorrelação serial (cerca de 0,73 na soja) e clara separação de nível antes e depois de 2006; segundo, abrir e destocar mata até o primeiro plantio é um investimento plurianual sujeito a expectativas.
    \item Por isso, delimitamos rigorosamente o alcance da inferência: o modelo não estima um experimento aleatório ano a ano, mas associações condicionais robustas que revelam como ciclos de mercado e instituições moldam a fronteira agrícola ao longo de duas décadas.
\end{enumerate}

\newpage

\renewcommand{\slidenum}{15}
% ============================================================
% SLIDE 15 / 15
% ============================================================
\noindent
\begin{minipage}{\textwidth}
\begin{center}
{\color{cinza}\footnotesize\textbf{TÓPICOS EM MICROECONOMIA APLICADA --- 16 DE NOVEMBRO DE 2026}\\[1.2mm]}
{\color{cinza}\bfseries\Large SLIDE 15 / 15\\[1.5mm]}
{\color{verde}\bfseries\LARGE Síntese da evolução e próximos passos da tese\\[2.5mm]}
\rule{\textwidth}{1pt}
\end{center}
\end{minipage}
\vspace{1.5mm}

\begin{tcolorbox}[title={\bfseries\small ÂNCORAS EMPÍRICAS E CONCEITUAIS DO SLIDE}]
\begin{itemize}[leftmargin=*,itemsep=1mm,topsep=0.5mm]
    \item \textbf{Avanço (OLS $\rightarrow$ FE):} H1a validada sob alto insumo; H1b validada para embargos ($-23{,}3\%$ no 1º termo), e não para multas avulsas; vazamento Amazônia--Cerrado e papel da ferrovia/pecuária.
    \item \textbf{Capítulo 2 (Próximo passo):} Via reversa na margem intensiva — efeitos da perda acumulada de capital natural e microclima sobre produtividade agrícola e resiliência da 2ª safra.
\end{itemize}
\end{tcolorbox}
\vspace{2mm}

\noindent{\color{verde}\large\bfseries Roteiro de Apresentação Oral}
\vspace{0.8mm}
\begin{enumerate}[leftmargin=*,label=\textbf{(\arabic*)},itemsep=2.8mm,topsep=1mm]
    \item Para concluir: a evolução desde as regressões OLS preliminares de 27 de outubro até a especificação em painel com efeitos fixos e interações resolveu os quebra-cabeças iniciais e consolidou o Capítulo 1 da tese.
    \item Mostramos que o mecanismo de mercado previsto em H1a opera claramente nos dados observados, mas condicionado ao pacote tecnológico de alto insumo e concentrado nas frentes de expansão e na dinâmica preparatória via pastagens.
    \item Evidenciamos que a dissuasão ambiental prevista em H1b funciona onde há incidência patrimonial e restrição de crédito — via embargos ambientais — e que a proteção setorial restrita à Amazônia conviveu com pressões persistentes sobre o Cerrado vizinho.
    \item Com esse primeiro elo entre incentivos econômicos, instituições e mudança de cobertura fechado, o próximo passo da tese é avançar para o Capítulo 2: investigar a via reversa na margem intensiva, analisando como a transformação da paisagem e o microclima regional afetam a produtividade e a resiliência das lavouras.
    \item Muito obrigado pela atenção, e fico aberto às perguntas e sugestões da turma e do professor.
\end{enumerate}

\newpage
\end{document}